\newpage
\chapter{秩序の維持、解任及び処分}

\section{相談又は問題を認知したとき}

担当者は、関係する投稿や申出の内容、時点、当事者の説明を確認する。
相談だけを望む者には希望を確かめ、急迫した危害がない限り、一方の説明だけで違反を確定しない。
公示済みのコミュニティルールを参照するときは、その本文と適用場面を確認する。
ネタバレへの配慮やロール選択などのお願いを、それだけで学園規則第８章第４条の処分事由に読み替えない。
処分を検討する場合は、同条のどの事由に該当するかを具体的に示す。

\section{通常の対応と処分}

問題を確認した担当者は、当事者への連絡や事実の訂正で対応できるか検討し、必要なら他の教務主事と共有する。
処分を行う場合は、本人の説明を確認し、学園規則第８章第４条の事実と第５条の処分の種類を分けて記録する。
警告、活動停止、除名の順を原則とし、警告後の同一学期内の再違反、活動停止中又は終了後の再違反は第６条に従って判断する。
警告を経ずに活動停止又は除名を行えるのは、第７条の著しく重大な違反がある場合に限る。
本人には、認定した事実、根拠、処分内容と期間がある場合はその期間、不服申立て先を伝える。
担当者は判断に参加した者、通知、開始・終了時点を記録する。

\section{緊急の措置}

秩序又は学生の安全への危害が現に生じ、又は急迫しているときは、学園規則第８章第８条により、必要な範囲で一時的なアクセス制限等を行える。
担当者は措置の範囲と開始時点を記録し、遅滞なく本人に理由を通知して、他の教務主事に共有する。
措置は当該学期の末日までに効力を失う。
これは処分ではないため、継続的な制限が必要なら別途、同章の処分の要否を判断する。

\section{辞任、解任及び不服申立て}

役職者の辞任の届出先は、学園長以外は学園長、学園長は教務主事である（第８章第２条）。
解任は第３条の事由を確認し、学園長以外の役職者は学園長が教務主事との協議を経て、学園長は教務主事の協議で判断する。
解任は学生としての処分ではなく、必要なら処分を別に検討する。
本人には解任の事実と理由を伝え、判断者及び職務の引継ぎを記録する。

解任、緊急の措置又は処分への不服申立てを受けたときは、第８章第９条に従い、元の判断を行った者以外の教務主事が審査する。
審査担当者は事実、適用条文及び措置の必要性を確認し、結論と理由を申立人へ伝えて記録する。
申請の不認可の再確認は、この不服申立てと区別して第３章で扱う。
